% Figure: rocket-nozzle thrust versus altitude for a nozzle sized for a
% fixed design altitude. Vacuum thrust is the ceiling; sea-level thrust
% is reduced by the ambient back-pressure term p_a A_e. The momentum
% thrust is constant; the pressure-thrust term grows with altitude as
% the ambient pressure falls.
\begin{figure}[htbp]
\centering
\begin{tikzpicture}
\begin{axis}[
  width=12cm, height=8cm,
  xlabel={altitude (km)}, ylabel={thrust (fraction of vacuum thrust)},
  xmin=0, xmax=60, ymin=0.78, ymax=1.02,
  axis lines=left,
  grid=major, grid style={enggray!20},
  tick label style={font=\scriptsize},
  label style={font=\small},
  legend style={font=\scriptsize, at={(0.97,0.03)}, anchor=south east},
  legend cell align=left,
]
% Thrust(h) = F_vac - p_a(h) * A_e. Model the momentum term + p_e*A_e as a
% constant C and subtract p_a(h)*A_e. Normalise so vacuum (p_a=0) gives 1.
% Use exponential atmosphere p_a = p0*exp(-h/H), H=7 km, and a coefficient
% k = p0*A_e / F_vac chosen so sea-level loss is about 0.18.
% F/F_vac = 1 - k*exp(-x/7), k=0.18.
\addplot[engblue, very thick, domain=0:60, samples=120] {1 - 0.18*exp(-x/7)};
\addlegendentry{delivered thrust $F(h)/F_{\mathrm{vac}}$}
% horizontal vacuum ceiling
\addplot[enggray, dashed, domain=0:60, samples=2] {1.0};
\addlegendentry{vacuum ceiling}
% design-altitude marker: where p_a = p_e, here taken as 12 km
\draw[engred, dashed] (axis cs:12,0.78) -- (axis cs:12,1.02);
\node[engred, font=\scriptsize, anchor=south, rotate=90] at (axis cs:12,0.86)
  {design altitude $p_a=p_e$};
% sea-level loss annotation
\draw[engamber, {Latex[length=2mm]}-{Latex[length=2mm]}]
  (axis cs:1.5,0.82) -- (axis cs:1.5,1.0);
\node[engamber, font=\scriptsize, anchor=west] at (axis cs:2.5,0.91)
  {pressure loss $p_a A_e$};
\end{axis}
\end{tikzpicture}
\caption[Rocket-nozzle thrust against altitude]{Delivered thrust of a
fixed-geometry rocket nozzle as a function of altitude, normalised to
the vacuum thrust. The thrust equation $F = \dot m\,u_e + (p_e - p_a)
A_e$ splits into a constant momentum-plus-exit-pressure part and a
back-pressure penalty $p_a A_e$ that shrinks as the ambient pressure
$p_a$ falls with altitude. At sea level the penalty is largest and the
nozzle delivers the least thrust; in vacuum the penalty vanishes and
the thrust reaches its ceiling. The nozzle is sized so that the exit
pressure matches the ambient pressure ($p_e = p_a$) at one design
altitude, marked by the dashed line; below it the nozzle is
overexpanded, above it underexpanded. The model uses an exponential
atmosphere with scale height $7\,\si{\kilo\meter}$ and a sea-level
penalty of about $18\,\%$, representative of a first-stage nozzle sized
for moderate altitude.}
\label{fig:vol03ch13:nozzle-thrust-altitude}
\end{figure}
